\section{Layers and cohomological tests}
\label{sec:layers}

We turn the twist action into a calculus of generator layers.  The
objectives are to detect the top and socle of a layer by ordinary
cohomology and to bound the layers created by a reduced spherical word.
These bounds will make the later cap-preserving operations possible.

We work in the homotopy category of bounded complexes of finite graded
left projective $Z$-modules.  A differential has cochain degree $1$ and
internal degree $0$.  The internal shift $\langle u\rangle$ raises
internal degrees by $u$, whereas $[1]$ lowers cochain placements by one.
Thus a generator of a copy of $P_\alpha\langle u\rangle$ in cochain
degree $c$ has \emph{level}
\begin{equation}
 d=c+u.
 \label{eq:layers-level}
\end{equation}
Generator dimensions always count copies of projectives, rather than
the vector-space dimensions of the projectives themselves.

\subsection{Minimal complexes and the linear heart}

A bounded projective complex is \emph{minimal} if each differential
entry belongs to the graded radical of the additive category of
projectives.  Any invertible entry between shifted indecomposables can
be removed by Gaussian elimination, splitting off a contractible
two-term complex.  There are finitely many summands, so this procedure
ends in a minimal complex.  Such a model is unique up to chain
isomorphism: a homotopy equivalence between two minimal complexes is
termwise invertible modulo the additive radical, since the homotopies
are composed with radical differentials.  Each term map is therefore
invertible, by the Krull--Schmidt property of finite graded module
terms.  This also proves the corresponding minimality statement for
any of the finite additive categories of module or bimodule terms
used below.

In a minimal complex a nonzero differential entry has path length
$a=1$ or $2$.  Its generator placements change by $(1,-a)$, and hence
its level changes by $1-a$.  The entries of length one preserve the
level, and those of length two lower it by one.  For a minimal complex
$E$, let $Y^d(E)$ be the complex consisting of the generators of level
$d$, with the length-one part of the differential.  This is a complex:
the length-two component of the equation $\partial^2=0$ is exactly
the square of its length-one part.

More concretely, write $V^d_{\alpha,c}(E)$ for the multiplicity space
of $P_\alpha\langle d-c\rangle$ in cochain degree $c$.  A coefficient
of a length-one differential is a map
\[
 V^d_{\alpha,c}(E)\longrightarrow V^d_{\beta,c+1}(E)
\]
for a chosen functional on $e_\alpha Z^1e_\beta$.  These coefficient
maps satisfy the backtrack relation obtained by multiplying opposite
arrows in $Z^2$.  There are no other length-two products to impose.
All these arrays are finite, and every coefficient arrow raises $c$.
In particular their paths act nilpotently.

At an unfolded vertex $\alpha$, the \emph{top} of such an array is
the generator space modulo the images of all incoming coefficient
maps, and its \emph{socle} is the common kernel of all outgoing
coefficient maps.  These spaces retain their cochain and internal
placements.  We write $\operatorname{top}_\alpha M$ and
$\operatorname{soc}_\alpha M$ for them, and write
$\dim M_\alpha$ for the sum of the generator multiplicities over
all placements at $\alpha$.

\begin{lemma}[Layer structure]\label{lem:layer-structure}
The full subcategories defined by minimal generator levels $d\leq0$
and $d\geq0$ form a bounded $t$-structure.  Its truncation
$\tau_{\leq k}E$ is represented by the subcomplex of a minimal model
on levels at most $k$; the complementary quotient represents
$\tau_{\geq k+1}E$.

The heart is equivalent to the abelian category of finite graded
coefficient arrays just described.  In particular, subobjects of a
pure layer, after shifting that layer into the heart, are precisely
graded subrepresentations of its array.  The unshifted cohomology
layer at $d$ is $Y^d(E)$.
\end{lemma}

\begin{proof}
The level calculation shows that the indicated low terms form a
subcomplex, with the indicated high quotient.  An internal- and
cochain-degree-zero map from a generator of level $d$ to one of
level $e$ would require a path of length $d-e$.  Consequently every
such map from levels at most $0$ to levels at least $1$ is zero.
This gives the required orthogonality.  The cochain shift $[1]$
decreases all levels by one, giving the two shift-closure axioms.
The short exact sequence of low subcomplex, full complex, and high
quotient gives the truncation triangle.  There are finitely many
occupied levels, so the $t$-structure is bounded.

Between pure objects of the same level, a degree-zero map has only
length-zero entries.  It is therefore exactly a graded map of their
multiplicity arrays commuting with the coefficient arrows.  A
cochain homotopy would have cochain degree $-1$ and would require
path length $-1$, so there are no homotopies to quotient out.
Conversely every finite array satisfying the backtrack relation
defines a pure complex of projectives.  Kernels and cokernels of
array maps, taken in each graded coefficient space, still satisfy
the relation.  This identifies the heart with the asserted abelian
category.  Applying the truncations twice identifies its cohomology
layers with the complexes $Y^d(E)$.
\end{proof}

We use the usual long exact cohomology sequence of a $t$-structure
\cite[Theorem~1.3.6]{BBD1982}.  To distinguish it from the vector-space tests
below, denote the heart-valued cohomology by $H^d_{\mathrm{lay}}$;
our conventions give
\[
 H^d_{\mathrm{lay}}(E)=Y^d(E)[d],\qquad
 H^d_{\mathrm{lay}}(E[1])=H^{d+1}_{\mathrm{lay}}(E).
\]
The corresponding fusion-category linear-heart construction appears in
\cite[\S6.1]{HengLicata2024}.

An exact functor has \emph{layer amplitude} $[a,b]$ if it sends each
pure layer at $d$ to levels in $[d+a,d+b]$.  By bounded truncations,
the same statement gives the corresponding bounds on an arbitrary
interval of input levels.

\begin{lemma}[Adjacent-layer exact sequences]
\label{lem:layers-amplitude-extension}
Suppose that an exact functor $F$ has layer amplitude $[0,1]$.
For any $E$ and any $d$, there is a short exact sequence
\begin{equation}
\begin{split}
0\longrightarrow Y^d\bigl(F(Y^{d-1}(E))\bigr)
\longrightarrow Y^d(FE)
\longrightarrow Y^d\bigl(F(Y^d(E))\bigr)
\longrightarrow0,
\end{split}
\label{eq:layers-adjacent-positive}
\end{equation}
where all terms are shifted by $[d]$ when regarded as objects of the
heart.  In particular their generator dimensions are additive.

For amplitude $[-1,0]$, the corresponding sequence is
\begin{equation}
0\longrightarrow Y^d\bigl(F(Y^d(E))\bigr)
\longrightarrow Y^d(FE)
\longrightarrow Y^d\bigl(F(Y^{d+1}(E))\bigr)
\longrightarrow0.
\label{eq:layers-adjacent-negative}
\end{equation}
\end{lemma}

\begin{proof}
In the positive case apply $F$ to
\[
 \tau_{\leq d-1}E\longrightarrow\tau_{\leq d}E
 \longrightarrow Y^d(E).
\]
The first term has levels at most $d$, and the last term has levels
$d,d+1$.  The flanking groups in the long exact sequence at $d$ are
$H^{d-1}_{\mathrm{lay}}F(Y^d(E))$ and
$H^{d+1}_{\mathrm{lay}}F(\tau_{\leq d-1}E)$; both are zero.
This gives a short exact sequence at that level.  Applying $F$ to
the truncation below $d-1$ identifies its left term with the one in
\eqref{eq:layers-adjacent-positive}, since levels at most $d-2$ map
to levels at most $d-1$.  Applying it to the truncation above $d$
identifies the middle term with $Y^d(FE)$, since levels at least
$d+1$ remain at least $d+1$.

For the negative case use
$\tau_{\leq d}E\to\tau_{\leq d+1}E\to Y^{d+1}(E)$.
After $F$ the first term has levels at most $d$, and the last has
levels $d,d+1$.  The same exact-sequence argument, now discarding
input levels below $d$ and above $d+1$, gives
\eqref{eq:layers-adjacent-negative}.
\end{proof}

\begin{lemma}[Composition at the bottom]
\label{lem:layers-bottom-composition}
Suppose that each of $F_1,\ldots,F_r$ is exact and has no negative
change of layer level.  Starting from a pure layer at $d$, the
level-$d$ part of $F_r\cdots F_1$ is obtained by retaining only the
level-$d$ part after each intermediate operation.  In particular,
any relation between such composed functors gives the corresponding
relation between their operations on the bottom layer.
\end{lemma}

\begin{proof}
After the first operation the quotient by its level-$d$ truncation
has levels at least $d+1$.  Every subsequent functor keeps that
quotient in levels at least $d+1$.  The long exact sequence therefore
identifies bottom cohomology before and after discarding it.
Repeat this argument at each step.
\end{proof}

\subsection{Extremal tests and single twists}

When vector-space cohomology retains both gradings, we write
$H^t(V)$ for the direct sum of its cohomology spaces whose cochain
degree plus internal degree is $t$.  Differentials raise this total
degree by one.  Only cochain degrees enter the signs.

\begin{lemma}[Extremal cohomological tests]
\label{lem:layers-extremal-tests}
For a pure layer $M$ at level $d$,
\begin{equation}
 H^d(e_\alpha M)\cong\operatorname{soc}_\alpha M,
 \qquad
 H^{d+2}(e_\alpha M)\cong
                  (\operatorname{top}_\alpha M)\langle2\rangle,
 \label{eq:layers-extremal-tests}
\end{equation}
as graded multiplicity spaces, with their corresponding placements.
The shift on the top is the internal placement of the token.
If a minimal complex $E$ has no generators above $d$ at $\alpha$
or its neighbors, its test in degree $d+2$ is still
$(\operatorname{top}_\alpha Y^d(E))\langle2\rangle$, and all
higher test cohomology there vanishes.
If it has no generators below $d$ at those vertices, its test in
degree $d$ is the socle of $Y^d(E)$ at $\alpha$.

A nonzero pure layer supported on a vertex set has a nonzero top
and a nonzero socle somewhere in that set.
\end{lemma}

\begin{proof}
For a pure layer at $d$, the paths in $e_\alpha M$ have lengths
$0,1,2$, and their total degrees are respectively $d,d+1,d+2$.
At the length-zero end, being a cycle means being killed by every
outgoing coefficient map.  At the token end, the incoming maps span
exactly the images of the incoming coefficient maps, because the
opposite-arrow pairing is nondegenerate.  This proves
\eqref{eq:layers-extremal-tests}.

For the upper assertion, every term of $e_\alpha E$ comes from a
generator at $\alpha$ or a neighbor.  Its actual total degree is its
generator level plus its path length, and is therefore at most
$d+2$.  At degree $d+2$ only the level-$d$ token terms occur.
Their incoming differential consists of the indicated linear maps:
any additional contribution would need a higher generator at the
same vertex or a neighbor.  The lower assertion follows in the same
way from the length-zero end.  Finally, the coefficient arrows
strictly raise cochain placement in a finite array.  Iterating them
cannot avoid a socle, and iterating incoming images cannot generate
a nonzero array with zero top everywhere.
\end{proof}

Let $F_i=B_{\sigma_i}(-)$ be the positive color twist from
Proposition~\ref{prop:twist-action}.  Write $r_i$ for the product of
the commuting coordinate reflections in its unfolded fiber.  On
vectors of generator multiplicities this changes only that fiber.

\begin{lemma}[Positive layers]\label{lem:positive-layer}
Each $F_i$ has layer amplitude $[0,1]$.  Its upper piece on a pure
layer $M$ at $d$ is supported on the fiber of $i$, has no arrows,
and has multiplicity $\dim\operatorname{top}_\alpha M$ at a vertex
$\alpha$ of that fiber.  Its bottom piece has unchanged
multiplicities at all other vertices, and at $\alpha$ in the fiber
it has multiplicity
\begin{equation}
 \dim Y^d(F_iM)_\alpha
 =\sum_\beta(2\id-\widetilde C)_{\alpha\beta}
       \dim M_\beta
       -\dim M_\alpha+\dim\operatorname{top}_\alpha M.
 \label{eq:layers-positive-dimension}
\end{equation}
The bottom piece has zero socle in color $i$.  If $M$ has zero
socle in every color of a set $J$ containing $i$, its bottom piece
after $F_i$ also has zero socle in all of $J$.

For any bounded left complex $E$, without a projectivity assumption
on its terms, there is a natural homotopy equivalence of test
complexes
\begin{equation}
 e_\alpha F_iE\simeq(e_\alpha E)\langle2\rangle[1]
 \qquad(\alpha\text{ of color }i).
 \label{eq:layers-positive-test-shift}
\end{equation}
There is the identical statement for right complexes tested on the
right.
\end{lemma}

\begin{proof}
For pure input at $d$, the simultaneous evaluation cone inserts
generators in color $i$ at levels $d+a-1$, for path lengths
$a=0,1,2$.  Those at $d-1$ are the length-zero evaluation copies.
Their maps to the old generators of color $i$ are identities modulo
the additive radical.  Cancelling these pairs removes all
generators below $d$.  None were inserted above $d+1$, proving the
amplitude bound.  Only the acted color can occur above $d$.

For one vertex $\alpha$ of this color put $A=e_\alpha E$.
All other vertices of the same color are orthogonal to it, so its
simultaneous test is
\begin{equation}
 \Cone\bigl(A\oplus A\langle2\rangle
       \xrightarrow{(\id,\theta_\alpha)}A\bigr).
 \label{eq:layers-token-cone}
\end{equation}
Projection onto the token source summand
$A\langle2\rangle[1]$ is a chain map with kernel the cone of
$\id_A$.  It is therefore a homotopy equivalence, proving
\eqref{eq:layers-positive-test-shift} for arbitrary input.
The same calculation works on the right.  For pure input, the
degree-$d+3$ test on the output is the old degree-$d+2$ test.
Lemma~\ref{lem:layers-extremal-tests} identifies it with the upper
generator multiplicities and with the old top, respectively.
There are no arrows among the vertices of the upper piece.

For completeness, signed generator counts are computed with
$(-1)^d$.  A contractible scalar pair has consecutive levels, so
its two signed contributions cancel.  In the cone for one vertex,
the inserted signed contribution from a generator at $\beta$ is
minus the alternating path count
$2\delta_{\alpha\beta}-(2\id-\widetilde C)_{\alpha\beta}$.
Thus the total signed vector transforms by $r_i$.  Since the
output has only levels $d,d+1$, subtracting its upper vector from
its bottom vector gives $r_i\dim M$.  The upper vector is the old
top.  At the acted fiber this gives
\eqref{eq:layers-positive-dimension}; outside it the multiplicities
are unchanged.

At a vertex of color $i$, the degree-$d$ test of $F_iM$ is zero
by \eqref{eq:layers-positive-test-shift}, so its bottom socle is
zero.  For a vertex $\beta$ of another color the evaluation
triangle for $e_\beta F_iM$ has first term a sum of copies of
$(e_\alpha M)\langle1\rangle$, with $\alpha$ in the acted fiber
adjacent to $\beta$.  Its cohomology at $d$ is zero, and its
cohomology at $d+1$ is a sum of the old socles at those $\alpha$.
If the socles in $J$ vanish, the long exact sequence at $d$,
together with the old socle vanishing at $\beta$, gives the
claimed vanishing at $\beta$ after the move.
\end{proof}

\begin{lemma}[Inverse amplitude]\label{lem:layers-inverse-amplitude}
If an exact equivalence $F$ has layer amplitude $[0,1]$, its inverse
has amplitude $[-1,0]$.  In particular each $F_i^{-1}$ has that
amplitude and satisfies
\begin{equation}
 e_\alpha F_i^{-1}E\simeq
       (e_\alpha E)\langle-2\rangle[-1]
 \qquad(\alpha\text{ of color }i),
 \label{eq:layers-negative-test-shift}
\end{equation}
with the analogous formula for right complexes.

For either sign, acting in color $i$ leaves the generator
multiplicities in every layer unchanged at all other colors.
\end{lemma}

\begin{proof}
It suffices to test the inverse on a heart object $X$.  If $A$ has
levels at most $-2$, then $FA$ has levels at most $-1$, whence
\[
 \Hom(A,F^{-1}X)=\Hom(FA,X)=0.
\]
Orthogonality implies that $F^{-1}X$ has levels at least $-1$.
If $B$ has levels at least $1$, then $FB$ does also, and
\[
 \Hom(F^{-1}X,B)=\Hom(X,FB)=0.
\]
It follows that $F^{-1}X$ has levels at most $0$.  This proves the
inverse bound; shifting proves it for every pure level.
Formula~\eqref{eq:layers-negative-test-shift} follows by applying
\eqref{eq:layers-positive-test-shift} to $F_i^{-1}E$ and using
invertibility.  This argument is at the homotopy level also for
nonprojective test inputs.

For $\beta$ outside the acted color, the right simple $U_\beta$
satisfies
\[
 U_\beta B_{\sigma_i}^{\pm1}\simeq U_\beta,
\]
because all the inserted projective bimodule terms have left
vertex in color $i$.  Tensoring a minimal left projective complex
with $U_\beta$ kills its radical differential and counts precisely
its generators at $\beta$, with both placements retained.  The
displayed identity therefore proves the last assertion in every
layer.
\end{proof}

\subsection{Positive divisors detected by cohomology}

Fix a spherical subset $J$.  Its unfolding is a finite simply-laced
Coxeter system by Lemma~\ref{lem:finite-unfolding}.  In the next two
lemmas the letters are its individual unfolded vertices.  For a
bounded left complex $X$ put
\begin{equation}
 h_J(X)=\max\{t:H^t(e_\alpha X)\ne0
                        \text{ for some }\alpha\in\widetilde J\},
 \label{eq:layers-test-height}
\end{equation}
and set $h_J(X)=-\infty$ if all these tests vanish.  For right
complexes replace $e_\alpha X$ by $Xe_\alpha$.  The statements
require only bounded test cohomology; the input need not be a
projective complex.

The next two lemmas use the extreme-degree method of
Brav--Thomas~\cite[Lemma~3.3 and Proposition~3.1]{BravThomas2011}.
They work with finite ADE configurations of $2$-spherical objects.
We give the total-degree calculation on both module sides, including
the nonprojective right simples needed below, and then pass from
unfolded vertices to the original colors.

\begin{lemma}[Change of test height]\label{lem:layers-height-change}
For a positive unfolded letter $\alpha$, if $h_J(X)$ is finite,
\[
 h_J(X)\leq h_J(B_\alpha X)\leq h_J(X)+1.
\]
An increase occurs exactly when $\alpha$ is an old highest-test
vertex; in that case it is the only new highest-test vertex.
Vanishing of all tests persists.  The same assertions hold for
right actions.
\end{lemma}

\begin{proof}
The test at $\alpha$ shifts up by one in total degree.  A
nonneighbor's test is unchanged.  At a neighbor $\beta$ in
$\widetilde J$ the evaluation triangle is
\begin{equation}
 (e_\alpha X)\langle1\rangle\longrightarrow e_\beta X
       \longrightarrow e_\beta B_\alpha X.
 \label{eq:layers-neighbor-triangle}
\end{equation}
If the old maximum is $t$, the neighboring new tests cannot
exceed $t$.  Thus the maximum rises precisely when the test at
$\alpha$ was nonzero at $t$, and only $\alpha$ attains $t+1$.
If its old maximum was $t-1$, its shifted test reaches $t$.
If it was below $t-1$, any nonzero old degree-$t$ test at another
vertex survives, by the long exact sequence of
\eqref{eq:layers-neighbor-triangle}.  This proves monotonicity.
The same triangles prove persistence of total vanishing.  Reversing
the module sides proves the right-action statements.
\end{proof}

\begin{lemma}[Top-divisor test]\label{lem:top-divisor}
Let $b$ be a positive word in the unfolded finite-type Artin
generators.  If
\[
 h_J(B_bX)>h_J(X),
\]
then every unfolded vertex attaining the highest test cohomology
of $B_bX$ is a left divisor of $b$ in the positive Artin monoid.
For a right action, every such vertex is a right divisor of $b$.
These assertions hold for arbitrary inputs with bounded test
cohomology, including complexes of right simples.
\end{lemma}

\begin{proof}
We prove the left statement by induction on the positive word
length, writing $b=\alpha b'$ and $Y=B_{b'}X$.  If this last
action raises height, its only new top vertex is $\alpha$, and
the assertion follows immediately.  Otherwise let
$h_J(Y)=h_J(B_\alpha Y)=t>h_J(X)$.  The top vertex $\alpha$,
if present, already left-divides $b$.  A top vertex $\beta\ne
\alpha$ not adjacent to $\alpha$ was a top of $Y$; by induction
it divides $b'$, and it commutes past $\alpha$.

It remains to treat an adjacent top vertex $\beta$.  Since height
did not rise, $H^t(e_\alpha Y)=0$.  The triangle
\eqref{eq:layers-neighbor-triangle} gives
\begin{equation}
 H^t(e_\beta B_\alpha Y)
 =\operatorname{coker}\bigl(
 H^{t-1}(e_\alpha Y)\longrightarrow H^t(e_\beta Y)\bigr).
 \label{eq:layers-divisor-cokernel}
\end{equation}
The nonzero cokernel implies that $\beta$ was a top of $Y$.
Induction allows $b'=\beta v$ using positive relations.  Put
$Z_0=B_vX$, so $Y\simeq B_\beta Z_0$.

We identify the arrow in
\eqref{eq:layers-divisor-cokernel} directly.  Set
$A=e_\beta Z_0$, $C=e_\alpha Z_0$, and choose arrows
\[
 r\in e_\alpha Z e_\beta,\qquad
 s\in e_\beta Z e_\alpha,\qquad
 sr=\kappa\theta_\beta,\quad\kappa\ne0.
\]
The two test complexes are
\[
 e_\beta Y=\Cone\bigl(A\oplus A\langle2\rangle
                   \xrightarrow{(\id,\theta_\beta)}A\bigr),
 \qquad
 e_\alpha Y=\Cone\bigl(A\langle1\rangle
                   \xrightarrow{r}C\bigr).
\]
Let $\pi:e_\beta Y\to A\langle2\rangle[1]$ be the token
projection of \eqref{eq:layers-token-cone}.  The arrow
$s:(e_\alpha Y)\langle1\rangle\to e_\beta Y$ maps its source
summand to $\kappa$ times the token source summand, and maps its
target summand into the target $A$, which $\pi$ kills.  Thus
$\pi s$ is exactly $\kappa$ times the connecting projection of
the second cone, up to a cochain sign.  This is an identity of
chain maps and needs no projectivity of $Z_0$.

Since $\pi$ is a homotopy equivalence, exactness of the second
cone identifies the image of the map in
\eqref{eq:layers-divisor-cokernel} with
\begin{equation}
 \ker\bigl(H^{t-1}(e_\beta Z_0)
                  \longrightarrow H^t(e_\alpha Z_0)\bigr).
 \label{eq:layers-divisor-kernel}
\end{equation}
The cokernel is nonzero, so the last map has nonzero image.
Consequently $H^t(e_\alpha Z_0)\ne0$.  Height monotonicity gives
$h_J(Z_0)\leq h_J(Y)=t$, hence $h_J(Z_0)=t>h_J(X)$.
Induction applied to $v$ makes $\alpha$ a left divisor of $v$.
Writing $v=\alpha v'$ and using the three-term braid gives
\[
 b=\alpha\beta\alpha v'=\beta\alpha\beta v',
\]
which extracts $\beta$ as required.  All cases of the induction
are proved.  The right-sided cones and the opposite arrows give
the right-divisor assertion by the same argument.
\end{proof}

A negative lift of a Coxeter expression is obtained by replacing each
letter by its inverse Artin generator, without reversing the order.

\begin{corollary}[Reduced-word amplitude]
\label{cor:reduced-amplitude}
For spherical $J$, a positive lift of any reduced expression in
$W_J$ has layer amplitude $[0,1]$.  A negative lift of a reduced
expression has amplitude $[-1,0]$.
\end{corollary}

\begin{proof}
First consider a reduced unfolded word.  Its test height can
increase at most once.  Indeed, at a second increase the new
letter was a highest-test vertex of the preceding result, whose
height already exceeded the starting one.  By
Lemma~\ref{lem:top-divisor} that letter could be extracted on
the last-action side of the preceding positive word.  Prepending
it would then give two equal consecutive Coxeter generators,
contrary to reducedness.

A pure input layer at $d$ has $h_J\leq d+2$.  Positives never
produce a layer below $d$.  Suppose that the output's highest
layer had level $a>d+1$.  Outside $\widetilde J$ generator
multiplicities in each layer are unchanged, so that highest
layer is supported in $\widetilde J$.  It has a nonzero top
there.  Lemma~\ref{lem:layers-extremal-tests} would give a test
at degree $a+2>d+3$, contrary to the preceding height bound.
If the input tests all vanish, they remain zero and give the
same contradiction.  This proves the positive bound in the
unfolding.

By Lemma~\ref{lem:finite-unfolding}, reduced expressions in
$W_J$ lift to reduced unfolded expressions, so the bound holds
for their positive color actions.  The inverse of a negative
reduced lift is a positive reduced lift of the reversed
expression.  Lemma~\ref{lem:layers-inverse-amplitude} therefore
gives the negative bound.
\end{proof}
